% ============================================================================
%   CHAPTER 3 — PROJECT CONTEXT AND SCOPE
%   Aivancity: 5-8 pages
%   ADAPTED for a PURELY ACADEMIC/RESEARCH context (no company internship)
% ============================================================================

\chapter{Project Context and Scope}
\label{chap:context}

% ----------------------------------------------------------------------------
\section{Academic Research Framework}
\label{sec:ctx-framework}

% [TODO — 1 page]
% - Aivancity's research environment
% - The PGE5 program structure
% - The relationship with the academic supervisor (Dr. Kar)
% - Nature of the research: pure academic, no industrial partner

% ----------------------------------------------------------------------------
\section{Project Origin and Evolution}
\label{sec:ctx-origin}

% [TODO — 1 page]
% - Initial topic: patient-reported uncertainty in trial matching
% - Reformulation after literature review: focus on LLM calibration
% - Justification: measurable, testable, addresses a documented gap

% ----------------------------------------------------------------------------
\section{Dataset: TrialGPT-Criterion-Annotations}
\label{sec:ctx-dataset}

\subsection{Dataset Description}
\label{subsec:ctx-desc}

% [TODO — 1.5 pages]
% - 1,015 patient-criterion pairs
% - Annotated by 3 physicians + GPT-4
% - 6-level ordinal eligibility labels
% - Provenance: Jin et al. 2024, TrialGPT (Nature Comm 15:9074)
% - Public availability, no patient data (synthetic notes)

\subsection{Justification of the Single-Dataset Choice}
\label{subsec:ctx-single}

% [TODO — 0.5 page]
% Explain the choice: TrialGPT-Criterion is the ONLY public dataset with
% - Expert (physician) annotations at the criterion-level (not trial-level)
% - Multi-level ordinal labels (6 levels)
% - Explicit criterion_type (inclusion vs exclusion)
% Other datasets (TREC CT, SIGIR, n2c2) either operate at trial-level or use
% only 3 levels, and thus cannot support fine-grained calibration analysis.

% ----------------------------------------------------------------------------
\section{Ethical Considerations}
\label{sec:ctx-ethics}

% [TODO — 0.5 page]
% - No real patient data (dataset uses synthetic notes)
% - Public data, CC BY 4.0
% - No IRB/CER approval required
% - Reproducibility commitments

% ----------------------------------------------------------------------------
\section{Constraints and Resources}
\label{sec:ctx-constraints}

\subsection{Computational Constraints}
\label{subsec:ctx-compute}

% [TODO — 0.5 page]
% - Kaggle free-tier T4 GPU (30h/week quota)
% - 4-bit quantization (bitsandbytes NF4) mandatory to fit 7-9B models
% - No 70B model runs feasible

\subsection{Model Access Constraints}
\label{subsec:ctx-access}

% [TODO — 0.5 page]
% - HuggingFace license approvals (Llama-3.1, Gemma)
% - Some models excluded due to licensing (e.g., some CC-BY-NC-only medical
%   LLMs)

\subsection{Time Constraints}
\label{subsec:ctx-time}

% [TODO — 0.5 page]
% - PFE conducted from March to September 2026 (6 months)
% - Design phase (March-April), experimental phase (May-July), writing phase
%   (August-September)

% ----------------------------------------------------------------------------
\section{Scope of This Work}
\label{sec:ctx-scope}

% [TODO — 1 page]
% What IS in scope:
%   - Systematic calibration audit on 12 LLM configurations
%   - Comparison of RLHF vs no-RLHF, generalist vs medical, prompt A vs B
%   - Post-hoc recalibration and conformal prediction
%   - Statistical analysis of pre-registered hypotheses
%
% What is NOT in scope (explicitly deferred):
%   - Multi-agent debate architectures (see future work)
%   - Fine-tuning or continued pretraining
%   - 70B models
%   - Multiple datasets validation
%   - Prospective clinical validation
